% ===== ch11 : 原文 main.tex 第 3658--4215 行（6b 等变 dg-范畴（含 dfn:GEquivQCohdg 起））=====
% 本文件是原文 LaTeX 源码的逐行分片，仅供翻译时对照；最终译文写到 src/chunks/ch11.tex。
\begin{dfn}\label{dfn:GEquivQCohdg}
% Let $X$ be a quasi-compact separated scheme over $\field$
% ---- 原文行 3660 ----
% satisfying the condition (GSP),
% and
% let $G$ be a group scheme acting on $X$.
% Let
% $\mathfrak U:=(U_i)_{i\in I}$
% be an affine open cover of $X$
% such that each $U_i$ is $G$-invariant.
The
\emph{$G$-equivariant coherent derived dg-category of $X$}
with respect to
$\mathfrak{U}$
is a dg-category
$$
\Coh_G(X;\mathfrak{U})_{\dg}
$$
defined in a similar way as Definition \ref{dfn:dg-enhancement-Coh}, where we replace $\check{C}^{\bullet} (\mathfrak{U},\mathcal{F}^{\bullet})$ by the $G$-invariant \v{C}ech cochain complex $\check{C}^{\bullet}_G (\mathfrak{U},(\mathcal{F}^{\bullet},\theta))$.
% as follows:

% \begin{itemize}
% \item
% ---- 原文行 3680 ----
% \emph{Objects} are bounded cochain complexes
% $$
% (\mathcal F^\bullet,\theta^\bullet)
% =
% \left(
% \bigoplus_{\objectdegree\in\mathbb Z}
% (\mathcal F^{\objectdegree},\theta^{\objectdegree}),
% d_{\mathcal F^\bullet}
% :=
% \left(d_{\mathcal{F}^{\bullet}}^i:\mathcal{F}^i\rightarrow\mathcal{F}^{i+1}\right)_{i\in\mathbb{Z}}
% \right)
% $$
% of finite-rank locally free $G$-equivariant
% $\mathcal O_X$-modules
% $(\mathcal F^{\objectdegree},\theta^{\objectdegree})
% \in\VB_G(X)$.
% Thus $(\mathcal F^{\objectdegree},\theta^{\objectdegree})=0$
% except for finitely many $\objectdegree$, and
% $$
% d_{\mathcal F^\bullet}^{\objectdegree}
% ---- 原文行 3700 ----
% \in
% \Hom_{\VB_G(X)}
% \left(
% (\mathcal F^{\objectdegree},\theta^{\objectdegree}),
% (\mathcal F^{\objectdegree+1},\theta^{\objectdegree+1})
% \right)
% \quad\text{satisfies }
% d_{\mathcal F^\bullet}^{\objectdegree}
% \comp
% d_{\mathcal F^\bullet}^{\objectdegree+1}
% =0
% \text{ for each $i$}.
% $$

% \item
% Given two objects
% $(\mathcal F^\bullet,\theta^\bullet)$ and
% $(\mathcal F'^\bullet,\theta'^\bullet)$,
% we first define their complex of $G$-equivariant internal homs
% $
% ---- 原文行 3720 ----
% (\mathcal H^\bullet,\Theta^\bullet)
% :=
% \mathcal{H}om_{\mathcal{O}_X}^{\bullet}
% \left(
% (\mathcal{F}^{\bullet},\theta^{\bullet}),
% (\mathcal{F}'^{\bullet},\theta'^{\bullet})
% \right)
% $
% by
% $$
% \mathcal{H}om_{\mathcal{O}_X}^{q}
% \left(
% (\mathcal{F}^{\bullet},\theta^{\bullet}),
% (\mathcal{F}'^{\bullet},\theta'^{\bullet})
% \right)
% :=
% \bigoplus_{\objectdegree\in\mathbb Z}
% \mathcal Hom_{\mathcal O_X}
% \left(
% (\mathcal F^{\objectdegree},\theta^{\objectdegree}),
% ---- 原文行 3740 ----
% (\mathcal F'^{\objectdegree+q},\theta'^{\objectdegree+q})
% \right)
% \quad\text{for }q\in\mathbb Z
% $$
% (see Definition \ref{defn:InternalHomOfEquivariantModules}),
% with a differential
% $$
% d_{\mathcal{H}^{\bullet}}^q
% :
% \mathcal{H}om_{\mathcal{O}_X}^{q}
% \left(
% (\mathcal{F}^{\bullet},\theta^{\bullet}),
% (\mathcal{F}'^{\bullet},\theta'^{\bullet})
% \right)
% \rightarrow
% \mathcal{H}om_{\mathcal{O}_X}^{q+1}
% \left(
% (\mathcal{F}^{\bullet},\theta^{\bullet}),
% (\mathcal{F}'^{\bullet},\theta'^{\bullet})
% \right)
% ---- 原文行 3760 ----
% $$
% defined by
% $$
% d_{\mathcal H^\bullet}f
% :=
% d_{\mathcal F^\bullet}\comp f
% -
% (-1)^q f\comp d_{\mathcal F'^\bullet}.
% $$
% Note that $(\mathcal{H}^{\bullet},\Theta^{\bullet})$ is a bounded cochain complex of
% finite-rank locally free $G$-equivariant $\mathcal{O}_X$-modules.
% \ \
% The \emph{dg-morphism space}
% from $(\mathcal{F}^{\bullet},\theta^{\bullet})$
% to $(\mathcal{F}'^{\bullet},\theta'^{\bullet})$
% is defined as the total complex of the bicomplex
% $\check{C}_G^\bullet
% \big(\mathfrak U,(\mathcal H^\bullet,\Theta^\bullet)\big)$,
% namely,
% it is a cochain complex of $\field$-modules whose degree-$k$ component is given by
% ---- 原文行 3780 ----
% \begin{align*}
% \hom_{\Perf_G(X;\mathfrak U)_{\dg}}^{\homdegree}
% \left(
% (\mathcal F^\bullet,\theta^\bullet),
% (\mathcal F'^\bullet,\theta'^\bullet)
% \right)
% &:=
% \bigoplus_{p+q=\homdegree}
% \check C_G^p
% \left(\mathfrak U,
% \mathcal{H}om_{\mathcal{O}_X}^{q}
% \left(
% (\mathcal{F}^{\bullet},\theta^{\bullet}),
% (\mathcal{F}'^{\bullet},\theta'^{\bullet})
% \right)
% \right)
% \\
% &
% \hspace{-30mm}
% =
% ---- 原文行 3800 ----
% \bigoplus_{p+q=\homdegree}
% \prod_{i_0<\cdots<i_p}
% \mathcal{H}om_{\mathcal{O}_X}^{q}
% \left(
% (\mathcal{F}^{\bullet},\theta^{\bullet}),
% (\mathcal{F}'^{\bullet},\theta'^{\bullet})
% \right)
% (U_{i_0\cdots i_p})^G
% \\
% &
% \hspace{-30mm}
% =
% \bigoplus_{p+q=\homdegree}
% \prod_{i_0<\cdots<i_p}
% \bigoplus_{\objectdegree\in\mathbb Z}
% \Hom_{\VB_G(U_{i_0\cdots i_p})}
% \left(
% (\mathcal F^{\objectdegree},\theta^{\objectdegree})
%   |_{U_{i_0\cdots i_p}},
% (\mathcal F'^{\objectdegree+q},\theta'^{\objectdegree+q})
% ---- 原文行 3820 ----
%   |_{U_{i_0\cdots i_p}}
% \right)
% \end{align*}
% (the last equality follows from Proposition \ref{prop:InvariantSectionsOfEquivariantInternalHom})
% and the \emph{differential} is given by the total differential
% $$
% D:=\delta+(-1)^p d_{\mathcal H^\bullet}.
% $$
% \item
% The \emph{composition}
% is defined by
% \begin{align*}
% (f\comp f')_{i_0\cdots i_{p+p'}}^{\objectdegree}
% :&=
% (-1)^{qp'}
% \left.
% f_{i_0\cdots i_p}^{\objectdegree}
% \right|_{U_{i_0\cdots i_{p+p'}}}
% \comp
% \left.
% ---- 原文行 3840 ----
% {f'}_{i_p\cdots i_{p+p'}}^{\objectdegree+q}
% \right|_{U_{i_0\cdots i_{p+p'}}}
% \\[3mm]
% &\in
% \Hom_{\VB_G(U_{i_0\cdots i_{p+p'}})}\left((\mathcal{F}^{\objectdegree},\theta^{\objectdegree})|_{U_{i_0\cdots i_{p+p'}}},(\mathcal{F}''^{\objectdegree+q+q'},\theta^{\objectdegree+q+q'}|_{U_{i_0\cdots i_{p+p'}}}\right)
% \end{align*}
% for
% \adjustbox{width=160mm}{
% $
% f
% :=
% \left(
% f_{i_0\cdots i_p}^{\objectdegree}
% \in
% \Hom_{\VB_G(U_{i_0\cdots i_p})}\left((\mathcal{F}^{\objectdegree},\theta^{\objectdegree})|_{U_{i_0\cdots i_p}},(\mathcal{F}'^{\objectdegree+q},\theta'^{\objectdegree+q})|_{U_{i_0\cdots i_p}}\right)
% \right)_{
% \substack{
% i\in\mathbb{Z}
% \\
% i_0<\cdots<i_p
% ---- 原文行 3860 ----
% }
% }
% \in
% \check{C}_G^p\big(
% \mathfrak{U},\mathcal{H}om_{\mathcal{O}_X}^q(
% (\mathcal{F}^{\bullet},\theta^{\bullet}),
% (\mathcal{F}'^{\bullet},\theta'^{\bullet})
% )
% \big),
% $
% }
% \adjustbox{width=160mm}{
% $
% f'
% :=
% \left(
% {f'}_{i_0\cdots i_{p'}}^{\objectdegree}
% \in
% \Hom_{\VB_G(U_{i_0\cdots i_{p'}})}\left((\mathcal{F}'^{\objectdegree},\theta'^{\objectdegree})|_{U_{i_0\cdots i_{p'}}},(\mathcal{F}''^{\objectdegree+q'},\theta''^{\objectdegree+q'})|_{U_{i_0\cdots i_{p'}}}\right)
% \right)_{
% ---- 原文行 3880 ----
% \substack{
% i\in\mathbb{Z}
% \\
% i_0<\cdots<i_p
% }
% }
% \in
% \check{C}_G^{p'}\big(
% \mathfrak{U},\mathcal{H}om_{\mathcal{O}_X}^{q'}(
% (\mathcal{F}'^{\bullet},\theta'^{\bullet}),
% (\mathcal{F}''^{\bullet},\theta''^{\bullet})
% )
% \big).
% $
% }
% \item
% The \emph{unit}
% $
% \id_{(\mathcal{F}^{\bullet},\theta^{\bullet})}
% :=
% ---- 原文行 3900 ----
% \left(\left(\id_{(\mathcal{F}^{\bullet},\theta^{\bullet})}\right)_{i_0\cdots i_p}^{\objectdegree}\right)_{
% \substack{
% i\in\mathbb{Z}
% \\
% i_0<\cdots<i_p
% }
% }
% \in
% \check{C}_G^p\big(
% \mathfrak{U},\mathcal{H}om_{\mathcal{O}_X}^q
% \left(
% (\mathcal{F}^{\bullet},\theta^{\bullet}),
% (\mathcal{F}^{\bullet},\theta^{\bullet})
% \right)
% \big)
% $
% of $(\mathcal F^{\bullet},\theta^{\bullet})$ is given by
% $$
% \left(\id_{(\mathcal{F}^{\bullet},\theta^{\bullet})}\right)_{i_0}^{\objectdegree}
% :=
% ---- 原文行 3920 ----
% \id_{(\mathcal{F}^{\objectdegree},\theta^{\objectdegree})|_{U_{i_0}}}
% \in
% \Hom_{\VB_G(U_{i_0})}\left((\mathcal{F}^{\objectdegree},\theta^{\objectdegree})|_{U_{i_0}},(\mathcal{F}^{\objectdegree},\theta^{\objectdegree})|_{U_{i_0}}\right)
% \quad\text{for each }
% i_0\in I,\
% \objectdegree\in\mathbb{Z}
% $$
% when
% $p=0$, $q=0$,
% and
% zero otherwise.
% \end{itemize}
% These operations satisfy $D^2=0$, associativity, the Leibniz rule,
% and the unit condition.
\end{dfn}
A direct verification shows that $\Coh_G (X;\mathfrak{U})_{dg}$ is a dg category, and is the dg enhancement of the derived category $D^b(\Coh_G (X))$ of quasi-coherent sheaves.
\begin{rmk}
    The existence of open affine cover of $X$ by $G$-invariants is a somewhat technical assumption. Due to the work of Sumihiro \cite{Sumihiro74} such cover exists when $X$ is normal and $G$ is a connected linear algebraic group, which is enough for the cases discussed in this paper.
\end{rmk}
Note that the strict $\widehat{G}$-action
% ---- 原文行 3940 ----
$
\phi:
\widehat{G}
\rightarrow
\Aut_{\str}(\Coh_G(X))
$
on $\Coh_G(X)$
naturally induces a strict $\widehat{G}$-action
$
\phi:
\widehat{G}
\rightarrow
\Aut_{\str}(\Coh_G(X;\mathfrak{U})_{\dg})
$
on
$\Coh_G(X;\mathfrak{U})_{\dg}$
(Remark \ref{rmk:EmbeddingOfCochainComplexes}).

For a closed subgroup scheme
$H\subseteq G$,
% ---- 原文行 3960 ----
the restriction functor
$\Res^G_H:\Coh_G(X)\rightarrow\Coh_H(X)$
(Definition \ref{defn:RestrictionOfEquivariantSheaves})
induces a dg-functor
$\Res_G^H:\Coh_G(X;\mathfrak{U})_{dg}\rightarrow\Coh_H(X;\mathfrak{U})_{dg}$,
%(Remark \ref{rmk:EmbeddingOfCochainComplexes}),
and
we have a restricted strict $H^\perp$-action
$
\phi:
H^\perp
\rightarrow
\Aut_{\str}(\Coh_G(X)_{\dg}).
$
\begin{defn}[Restriction dg-functor]
For a closed subgroup scheme
$H\subseteq G$,
recall the restriction functor
$\Res^G_H:\Coh_G(X)\rightarrow\Coh_H(X)$
(Definition \ref{defn:RestrictionOfEquivariantSheaves}).
% ---- 原文行 3980 ----
Suppose that there exists an affine open cover
$\mathfrak U:=(U_i)_{i\in I}$
of $X$
such that each $U_i$ is $G$-invariant.
Then each $U_i$ is automatically $H$-invariant.
We define the restriction dg-functor
between dg-categories
$$
\Res_G^H
:
\Coh_G(X;\mathfrak{U})_{\dg}
\rightarrow
\Coh_H(X;\mathfrak{U})_{\dg}
$$
whose mapping on objects and morphisms is given as follows:
\begin{itemize}
\item
Objects:
$
(\mathcal{F}^{\bullet},\theta^{\bullet})
% ---- 原文行 4000 ----
\mapsto
(\mathcal{F}^{\bullet},\theta^{\bullet}|_H)
$
where
$$
(\mathcal{F}^{\bullet},\theta^{\bullet}|_H)
:=
\left(
\bigoplus_{\objectdegree\in\mathbb Z}
(\mathcal F^{\objectdegree},\theta^{\objectdegree}|_H),
d_{\mathcal F^\bullet}
:=
\left(d_{\mathcal{F}^{\bullet}}^i:\mathcal{F}^i\rightarrow\mathcal{F}^{i+1}\right)_{i\in\mathbb{Z}}
\right).
$$
Here,
each
$(\mathcal{F}^i,\theta^i|_H)$
is the image of
$(\mathcal{F}^i,\theta^i)$
% ---- 原文行 4020 ----
under the restriction functor
$\Res^G_H:\Coh_G(X)\rightarrow\Coh_H(X)$,
and
$
d_{\mathcal{F}^{\bullet}}^i
\in
\Hom_{\Coh_H(X)}(\mathcal{F}^i,\mathcal{F}^{i+1})
$
is the image of the morphism
$
d_{\mathcal{F}^{\bullet}}^i
\in
\Hom_{\Coh_G(X)}(\mathcal{F}^i,\mathcal{F}^{i+1})
$
under the functor.
%Note that
%the $H$-equivariance of
%$d_{\mathcal{F}^{\bullet}}^i$
%follows automatically from
%its $G$-equivariance for each $i\in\mathbb{Z}$.
% ---- 原文行 4040 ----
\item
Morphisms:
Recall the definition of the hom space
\[
\adjustbox{width=170mm}{%
  $\displaystyle
  \hom_{\Coh_G(X;\mathfrak U)_{\dg}}^{\homdegree}
  \left(
    (\mathcal F^\bullet,\theta^\bullet),
    (\mathcal F'^\bullet,\theta'^\bullet)
  \right)
  =
  \bigoplus_{p+q=\homdegree}
  \prod_{i_0<\cdots<i_p}
  \bigoplus_{\objectdegree\in\mathbb Z}
  \Hom_{\Coh_G(U_{i_0\cdots i_p})}
  \left(
    (\mathcal F^{\objectdegree},\theta^{\objectdegree})
      |_{U_{i_0\cdots i_p}},
    (\mathcal F'^{\objectdegree+q},\theta'^{\objectdegree+q})
% ---- 原文行 4060 ----
      |_{U_{i_0\cdots i_p}}
  \right).$
}
\]
Each summand has a natural restriction
\[
\adjustbox{width=170mm}{$
\displaystyle
\Hom_{\Coh_G(U_{i_0\cdots i_p})}
\left(
(\mathcal F^{\objectdegree},\theta^{\objectdegree})
  |_{U_{i_0\cdots i_p}},
(\mathcal F'^{\objectdegree+q},\theta'^{\objectdegree+q})
  |_{U_{i_0\cdots i_p}}
\right)
\rightarrow
\Hom_{\Coh_H(U_{i_0\cdots i_p})}
\left(
(\mathcal F^{\objectdegree},\theta^{\objectdegree}|_H)
  |_{U_{i_0\cdots i_p}},
% ---- 原文行 4080 ----
(\mathcal F'^{\objectdegree+q},\theta'^{\objectdegree+q}|_H)
  |_{U_{i_0\cdots i_p}}
\right),
$}
\]
%since a $G$-equivariant map is also a $H$-equivariant map.
given by the functor
$\Res^G_H:\Coh_G(U_{i_0\cdots i_p})\rightarrow\Coh_H(U_{i_0\cdots i_p})$.
Therefore,
there is a natural restriction
$$
\hom_{\Coh_G(X;\mathfrak U)_{\dg}}^{\homdegree}
\left(
(\mathcal F^\bullet,\theta^\bullet),
(\mathcal F'^\bullet,\theta'^\bullet)
\right)
\rightarrow
\hom_{\Coh_H(X;\mathfrak U)_{\dg}}^{\homdegree}
\left(
(\mathcal F^\bullet,\theta^\bullet),
% ---- 原文行 4100 ----
(\mathcal F'^\bullet,\theta'^\bullet)
\right).
$$
\end{itemize}
\end{defn}

\begin{crl}
\label{crl:RestrictionFunctionEm}
Let $G$ be a diagonalizable group scheme acting on a scheme $X$,
and $H \subseteq G$ a closed subgroup scheme.
Let
$\mathfrak U:=(U_i)_{i\in I}$
be an affine open cover of $X$
such that each $U_i$ is $G$-invariant
(and hence $H$-invariant).\par
Then the restriction dg-functor
$
\Res_H^G:\Coh_G(X;\mathfrak{U})_{\dg}
\rightarrow
\Coh_H(X;\mathfrak{U})_{\dg}
% ---- 原文行 4120 ----
$
induces a quasi-embedding of the orbit category:
$$
\Coh_G(X;\mathfrak{U})_{\dg}/H^\perp
\hookrightarrow
\Coh_H(X;\mathfrak{U})_{\dg}.
$$
\end{crl}

\begin{proof}
Let $(\mathcal{F},\theta)$ and $(\mathcal{F}',\theta')$
be two objects in $\Coh_G(X;\mathfrak{U})_{\dg}$.
$$
\hspace{-80mm}
\hom_{\Coh_H(X)_{\dg}}
\left((\mathcal{F}^{\bullet},\theta^{\bullet}|_H),(\mathcal{F}'^{\bullet},\theta'^{\bullet}|_H)\right)
$$
$$
=
\displaystyle
% ---- 原文行 4140 ----
\bigoplus_{p+q=\homdegree}
\prod_{i_0<\cdots<i_p}
\bigoplus_{\objectdegree\in\mathbb Z}
\Hom_{\Coh_H(U_{i_0\cdots i_p})}
\left(
(\mathcal F^{\objectdegree},\theta^{\objectdegree}|_H)
  |_{U_{i_0\cdots i_p}},
(\mathcal F'^{\objectdegree+q},\theta'^{\objectdegree+q}|_H)
  |_{U_{i_0\cdots i_p}}
\right)
$$
By Proposition \ref{prop:HInvariantsOfHomSpace},
each summand is decomposed as
$$
\hspace{-30mm}
\Hom_{\Coh_H(U_{i_0\cdots i_p})}
\left(
(\mathcal F^{\objectdegree},\theta^{\objectdegree}|_H)
  |_{U_{i_0\cdots i_p}},
(\mathcal F'^{\objectdegree+q},\theta'^{\objectdegree+q}|_H)
% ---- 原文行 4160 ----
  |_{U_{i_0\cdots i_p}}
\right)
$$
$$
=
\bigoplus_{\lambda\in H^{\perp}}
\Hom_{\Coh_G(U_{i_0\cdots i_p})}
\left(
(\mathcal F^{\objectdegree},\theta^{\objectdegree})
  |_{U_{i_0\cdots i_p}}(\lambda),
(\mathcal F'^{\objectdegree+q},\theta'^{\objectdegree+q}|_H)
  |_{U_{i_0\cdots i_p}}
\right)
$$
On the other hand,
by definition of orbit categories of dg-categories,
$$
\hspace{-80mm}
\hom_{\Coh_G(X)_{\dg}/H^{\perp}}
\left((\mathcal{F}^{\bullet},\theta^{\bullet}|_H),(\mathcal{F}'^{\bullet},\theta'^{\bullet}|_H)\right)
% ---- 原文行 4180 ----
$$
$$
\hspace{-58mm}
=
\displaystyle
\bigoplus_{\lambda\in H^{\perp}}
\hom_{\Coh_G(X)_{\dg}}
\left((\mathcal{F}^{\bullet},\theta^{\bullet}|_H)(\lambda),(\mathcal{F}'^{\bullet},\theta'^{\bullet}|_H)\right)
$$
$$
=
\displaystyle
\bigoplus_{\lambda\in H^{\perp}}
\bigoplus_{p+q=\homdegree}
\prod_{i_0<\cdots<i_p}
\bigoplus_{\objectdegree\in\mathbb Z}
\Hom_{\Coh_G(U_{i_0\cdots i_p})}
\left(
(\mathcal F^{\objectdegree},\theta^{\objectdegree})
  |_{U_{i_0\cdots i_p}}(\lambda),
% ---- 原文行 4200 ----
(\mathcal F'^{\objectdegree+q},\theta'^{\objectdegree+q}|_H)
  |_{U_{i_0\cdots i_p}}
\right).
$$
Therefore,
we have
$$
\hom_{\Coh_H(X)_{\dg}}
\left((\mathcal{F}^{\bullet},\theta^{\bullet}|_H),(\mathcal{F}'^{\bullet},\theta'^{\bullet}|_H)\right)
=
\hom_{\Coh_G(X)_{\dg}/H^{\perp}}
\left((\mathcal{F}^{\bullet},\theta^{\bullet}|_H),(\mathcal{F}'^{\bullet},\theta'^{\bullet}|_H)\right)
$$
as cochain complexes of $\field$-vector spaces.
\end{proof}

